\FifthArticle{5}
From the order of the positive and negative values of $T$ and $U$ discussed in the preceding article, which holds for every value of $r$ greater than $R$\SourceLeaf{086}\HistoricalNote{14}, it can now be inferred that within the domain of smaller values of $r$ there must be certain crossings in these arrangements which contain the essence of the theorem we are to prove. I shall present the proof in a form drawn from the geometry of position, because this gives it the greatest clarity and simplicity. At bottom, however, the actual content of the whole argument belongs to a higher domain of the general abstract theory of magnitudes, independent of spatial considerations, whose subject is combinations of magnitudes connected through continuity. This domain has so far been little cultivated, and one cannot move within it without a language borrowed from spatial images.

\FifthArticle{6}
The entire domain of complex quantities is represented by an unbounded plane, in which each point whose coordinates relative to two mutually perpendicular axes are $t,u$ is regarded as corresponding to the complex quantity $x=t+iu$. If this complex quantity is put in the form $x=r(\cos\rho+i\sin\rho)$, then $r,\rho$ denote the polar coordinates of the corresponding point. The collection of all complex quantities for which $r$ has one fixed value is therefore represented by a circle whose radius is this value and whose centre is the origin of coordinates. I shall denote by $K$ that circle for which $r$ exceeds $R$ by a difference chosen at will, and by $(1),\,(2),\,(3)\,.\,.\,.\,.\,(2n)$ those points on it corresponding to the values of $\rho$ lying respectively between $\omega$ and $3\omega$, between $5\omega$ and $7\omega$, between $9\omega$ and $11\omega$, and so on, up to between $(8n-3)\omega$ and $(8n-1)\omega$, for which, by Article 4, $T=0$. Notice that $U$ will be positive at the points $(1),\,(3),\,(5)$, etc., and negative at the points $(2),\,(4),\,(6)$, etc.
\SourceLeaf{087}
\FifthArticle{7}
The collection of points in our plane for which $T$ is positive forms connected regions, as is evident on considering that, in a continuous passage from one point to another, $T$ changes continuously. Likewise, all the points for which $T$ is negative form connected regions. Between the regions of the first kind and those of the second lie points at which $T=0$; by the nature of the function $T$, these points cannot themselves form areas, but only lines bounding regions of the one kind on one side and regions of the other kind on the other.

The space outside $K$ contains $n$ regions of the first kind alternating with just as many of the second kind, each extending continuously to infinity from a portion of the circumference $K$. At the same time, it is clear that each of these regions continues beyond the circumference into the interior, and that the following cases can occur as regards its further shape.

1) The region in question, beginning at a portion of $K$, ends in isolation within the disk. Its peripheral boundary then consists of only two connected pieces, one of which is part of $K$, while the other lies within the circle. In the accompanying figure, which concerns an equation of the fifth degree and in which the signs of $T$ are written in the various regions, three of the regions with positive $T$ fall into this case. One has the boundary lines $10.1$ and $1.11.10$, the second has $4.5$ and $5.12.4$, and the third $6.7$ and $7.13.6$. There are two similar regions with negative $T$.

2) The region simply crosses the disk so as to form one connected region with another entering at a different place. The entire peripheral boundary will then consist of four pieces, two belonging to the circumference $K$ and the other two to the interior. In our figure this case occurs in the region bounded by $2.3;$ $3.0.8;\ 8.9;\ 9.11.2$.

3) The region branches once or several times within the circle, so as to form one connected region with two or more others entering at other places. Its entire peripheral \SourceLeaf{088}boundary will then consist of six, eight, or more pieces in an even number, belonging alternately to the circumference and to the interior. In our figure this occurs in a region whose boundary is formed by the six pieces $3.4;\ 4.12.5;\ 5.6;\ 6.13.7;\ 7.8;\ 8.0.3$, but in which $T$ is negative.

\FifthArticle{8}

A complete enumeration of all conceivable shapes of the regions entering the circle would require further modifications of the cases given. For example, if such a region divides into two branches which join again in the interior, this could be included in the first or second case of the preceding article, according as, after they join, the region ends in the interior or continues (without dividing again) to another place on the circumference. Its shape would be modified only by enclosing an island that does not belong to it. Moreover, it would not be difficult to prove rigorously that, owing to the special nature of the function $T$, modifications of this kind are \emph{not possible at all}. For our purpose, however, this is unnecessary, since only the succession of pieces of the \emph{outer} boundary of each region in question matters (that is, of those in which $T$ is positive).

We have already observed that the number of these pieces is always even (two in the first case of the preceding article, four in the second, six or more in the third), one belonging to the circumference $K$ and one to the interior, alternately. Further, it is clear that if this outer boundary is traversed always in the same direction, here chosen as that in which the numbers of the points of $K$ increase (thus, for example, in our figure, so that the region always lies to the right of the boundary), the initial and final points of a piece belonging to the circumference will be designated respectively by an even number and the odd number greater by one. Thus the initial and final points of every piece passing through the interior are always designated respectively by an odd and an even number.

It is therefore established that each of the $n$ lines which enter the interior at a point of $K$ designated by an odd number\SourceLeaf{089}, and on which $T=0$ everywhere, passes continuously through this space in a completely definite way\HistoricalNote{15}, until it emerges at another place designated by an even number. Now, as already observed above (at the end of Article 6), the value of $U$ is positive at its initial point and negative at its final point. By continuity of the change in value, $U=0$ must therefore hold at an intermediate point. This point then represents a root of the equation $X=0$; and since the number of such lines $=n$, $n$ roots of that equation always result in this way.

\FifthArticle{9}

If the lines in question pass through the circle without meeting one another, it is clear that the $n$ roots obtained in this way are necessarily distinct. Such an unobstructed passage occurs in our figure for the lines from $3$ to $8$, from $5$ to $4$, and from $7$ to $6$, to which belong the roots represented by the points $0,\ 12,\ 13$. If, however, two or more such lines have a point in common, it is not thereby necessary, but is nevertheless possible, that this point is also the one at which $U=0$. In that case two or more roots coincide in one, or, as is usually expressed, are equal to one another. In our figure the lines $1.10$ and $9.2$ meet at the point $11$, and at this point $U=0$ also holds. Thus, besides the three distinct roots already listed, the equation has two equal roots.

